% Drawing 80/03-1 -- folder 80, batch 03, archive page 44 (PDF page 45).
% Transcription: à gauche du texte, un dessin : un heptagone aux côtés fléchés, sommets en pointillé, avec « $u_0(i) = i+1$ » et une flèche courbe au centre. « $d_i$ » est ajouté au-dessus de « bijection », « de $\vec{I}_i$ » est \uncertain{lu}. Au-dessus de « circulaire, i.e. transitive », une insertion : « (\ill{} par $i_\alpha = u_0^\alpha(i_0)$) ». « $\nu = \nu_s$ » est écrit dans la marge gauche, en face de (28)
% Kind: heptagon with arrowed sides. Method 2 (schematic redraw from the facsimile). Status: partial.
% Notes: Six sides drawn as arrows going round counter-clockwise on the page: top to upper left, down the left, along the bottom, up to the right corner, up to the upper right corner, back to the top. The six corners are computed once as the meeting points of consecutive sides, and each side lies on the line between its two corners, cut short where he left the corner open (the top, the lower left and lower right corners). The dotted points (mostly in pairs) outside the corners are placed as read, with the dash above the top corner; i under the top corner and u_0(i) = i+1 at \small (Michel's size), the curved arrow at the centre. Still differs: six arrowed sides and six corners are on the page, where the transcription says a heptagon; the seventh vertex is presumably implied by the dotted points, which are not joined. Redrawn 2026-11-09 (issue 128): to check again.
% Coordinates: 1 unit = 1 mm = 1/100 of the crop width (bbox in page fractions: [0.12, 0.3, 0.42, 0.45]).
% Redrawn by Claude (claude-opus-5-5), issue #73, 2026-11-09; not checked by a human; not compiled (no TeX in the session).
% Revised 2026-11-07 after Michel's check: labels enlarged.
% Rights: reproduces part of the work -- behind RENDER_TRANSCRIPTS, never published by itself.
\documentclass[tikz,border=3pt]{standalone}
\usepackage{amsmath,amssymb}
\usetikzlibrary{patterns,arrows.meta}
\tikzset{pencil/.style={black!45}, struck/.style={black!35}, mark/.style={circle,draw,inner sep=1.3pt}, pt/.style={circle,fill,inner sep=0.9pt}, opt/.style={circle,draw,fill=white,inner sep=1.1pt}}
\begin{document}
\begin{tikzpicture}[x=1mm,y=1mm,line cap=round,line join=round,font=\scriptsize]
\draw[-{Stealth[length=2mm]}] (50.7,-4.3) -- (30.2,-20.6);
\draw[-{Stealth[length=2mm]}] (30.4,-22.6) -- (33.4,-49.1);
\draw[-{Stealth[length=2mm]}] (35.5,-52.1) -- (72,-52.1);
\draw[-{Stealth[length=2mm]}] (75.1,-51.7) -- (83,-36.9);
\draw[-{Stealth[length=2mm]}] (83,-35.2) -- (79.7,-19.9);
\draw[-{Stealth[length=2mm]}] (79.7,-19.9) -- (55.6,-4.6);
\node[pt] at (19.4,-18.3) {};
\node[pt] at (22.9,-20.1) {};
\node[pt] at (83.9,-18.2) {};
\node[pt] at (88.5,-15) {};
\node[pt] at (83.7,-34.1) {};
\node[pt] at (88.9,-37.7) {};
\node[pt] at (94.4,-38.8) {};
\node[pt] at (74.9,-58) {};
\node[pt] at (77.2,-62.7) {};
\node[pt] at (28.8,-59.4) {};
\node[pt] at (25.9,-61.9) {};
\draw (52.1,-3) -- (53.3,-2.9);
\node[inner sep=1pt] at (52.4,-8.8) {\small $i$};
\node[inner sep=1pt] at (41.5,-21.6) {\small $u_0(i)=i+1$};
\draw[-{Stealth[length=2mm]}] plot[smooth, tension=0.55] coordinates {(62.7,-36) (64.9,-30.2) (63.3,-25.7) (59.5,-23.3) (55,-23.8) (51.1,-26)};
\end{tikzpicture}
\end{document}
